\section{An exact positive finite-budget inference state}\label{sec:filter}
\subsection{Cubic Bernstein coefficients of the mechanical kernel}
Set $t=(R-l)/(u-l)$ and
$B_i^d(t)=\binom di t^i(1-t)^{d-i}$. Write $\Delta=u-l$.
The degree-three Bernstein coefficients of $R,R^2,R^3$ are respectively
\begin{equation}
 r_j=l+\Delta j/3,\quad
 s_j=l^2+2l\Delta j/3+\Delta^2j(j-1)/6,\quad
 d_j=l^{3-j}u^j,\qquad 0\le j\le3.
 \label{eq:elementary-coefficients}
\end{equation}
These formulas follow by expanding $R=l+\Delta t$ and using the Bernstein
representations of $1,t,t^2,t^3$. Therefore the raw numerator coefficients are
\begin{equation}
 h_j(a,\partial)=\pi s_j,\quad
 h_j(a,0,y)=a_0-\pi s_j-br_j,\quad
 h_j(a,1,y)=b\{r_j+\epsilon\cos(y-\theta)d_j\}.
 \label{eq:raw-coefficients}
\end{equation}
They are strictly positive, uniformly over modes and reports. Indeed
$l\le r_j\le u$, $l^2\le s_j\le u^2$, and $d_j\le u^2r_j$.
For the last inequality it suffices to check $j=0,1,2,3$: the middle cases
reduce to $3l^2\le u(2l+u)$ and $3l\le l+2u$.
The same $h_*$ used in \eqref{eq:lower-bound} is a coefficient lower bound.
The failure coefficients
\begin{equation}
 h_j(a,m,\dagger)=\int(1-g(x))h_j(a,x)\,\nu(dx)
 \label{eq:failure-coefficients}
\end{equation}
are positive and depend only on the four moments $m$. The sum of the
integrals of the raw coefficients is $a_0$ at every $j$, since the raw
polynomial normalizes identically. Thus failure coefficients are in
$[\eta a_0,(1-\eta)a_0]$. Accepted factors $g(x)$ are positive scalars
independent of $R$ and may be omitted from a likelihood state.

\begin{theorem}[Positive polynomial filter with a full parameter likelihood]\label{thm:filter}
After $n$ consumed cartridges, the radius likelihood, up to a positive
radius-independent scalar, is
\[
 P_c(R)=\sum_{i=0}^{3n}c_iB_i^{3n}(t),\qquad c_i>0,\quad\sum_i c_i=1.
\]
The initial state is $c=(1)$. On an accepted $x$, multiply by $H_a(R,x)$;
on $\dagger$, multiply by \eqref{eq:failure-poly}. For a current degree $d$
and cubic coefficients $h_j$, the unnormalized update is
\begin{equation}
 \widehat c_k=\sum_{i+j=k}
 \frac{\binom di\binom3j}{\binom{d+3}k}c_i h_j,
 \qquad c'_k=\widehat c_k/\sum_\ell\widehat c_\ell.
 \label{eq:bern-product}
\end{equation}
For a fixed prior $\pi_0$, the exact posterior is
\begin{equation}
 \pi_c(dR)=\frac{P_c(R)\pi_0(dR)}{\int P_c\,d\pi_0}.
 \label{eq:posterior}
\end{equation}
For the uniform prior on $I$, its density in $t$ is the beta mixture
\begin{equation}
 \sum_{i=0}^d c_i\,\mathrm{Beta}(i+1,d-i+1)(dt).
 \label{eq:beta-mixture}
\end{equation}
Given the one-step coefficient vector, the update costs $O(n)$ positive
arithmetic operations and stores $3n+1$ real coefficients. Computing the
gate moments or preprocessing an arbitrary prior is not included in this
coefficient-update count. It retains the whole likelihood curve on $I$ and all of
its finite derivatives, not just information at one radius.
\end{theorem}
\begin{proof}
For a realized history the chosen commands are fixed functions of that
history, not of the unknown radius. Its density is a product of the
one-cartridge densities from Theorem~\ref{thm:attainable}. An accepted
observation contributes $g(x)H_a(R,x)/a_0$, and a censored observation
contributes $H_\dagger(R;m)/a_0$. Discarding only the common positive scalars
gives the asserted polynomial. The identity
$B_i^dB_j^3=\binom di\binom3j\binom{d+3}{i+j}^{-1}B_{i+j}^{d+3}$ proves
\eqref{eq:bern-product}; every summand is nonnegative and each output
coefficient has a positive summand. Bayes' formula proves
\eqref{eq:posterior}. Since $\int_0^1 B_i^d=1/(d+1)$ and
$(d+1)B_i^d$ is a beta density, normalized coefficients give
\eqref{eq:beta-mixture}. There are only four summands per output coefficient.
\end{proof}

For example, under the uniform prior,
\[
 \E_{\pi_c}t^q=\sum_{i=0}^d c_i
 \frac{(i+1)(i+2)\cdots(i+q)}{(d+2)(d+3)\cdots(d+q+1)}.
\]
Moments up to order three determine a raw next-report density. This does
not mean that only those moments determine all future Bayesian decisions:
the full coefficient vector is retained for subsequent updates.

\subsection{Compact state strata}
Let $\mathcal H$ be a compact box containing all raw and failure cubic
coefficient vectors, with a strictly positive lower coordinate bound.
For each $n\le N$, take $\cS_n$ to be the image of $\mathcal H^n$ under
polynomial multiplication and coefficient normalization. It is compact and
all its coefficient coordinates have a positive lower bound depending on
$n$. This follows because a positive continuous function on a compact set
has a positive minimum, and every product coefficient is positive.
The actual reachable states form a subset; this possibly larger compact
state space is closed under every actual update. On $\cS_n$ the posterior
map is continuous in variation: coefficient convergence gives uniform
polynomial convergence, and the common positive lower bound controls the
normalizer for any fixed prior. Raw and failure updates are jointly
continuous on their respective compact domains. This directly supplies the
continuity and evidence properties used below.

\subsection{A general positive-polynomial closure theorem}
\begin{theorem}[Uniform degree elevation and rational-experiment closure]\label{thm:general}
Suppose a finite-reference-measure experiment has densities
$H_a(R,x)/D(R)$ on a compact interval, where $D>0$ is fixed, $H$ is a
polynomial of degree at most $q$, its power coefficients are uniformly
bounded, and $H_a(R,x)\ge h>0$ uniformly. Then there is a common finite
Bernstein degree $m$ at which every $H_a(\cdot,x)$ has strictly positive
coefficients bounded above and below uniformly. Positive bounded
radius-independent gates and their failure atoms preserve this property.
After $n$ reports the likelihood has the form
$P_n(R)/D(R)^n$, with a positive Bernstein numerator of degree at most $mn$.
Consequently a normalized $mn+1$ coefficient vector and $n$ retain the
entire likelihood and give exact positive updates and posterior prediction.
\end{theorem}
\begin{proof}
Transform the parameter interval to $[0,1]$ and write
$H(t)=\sum_{j=0}^q a_jt^j$. At degree $m\ge q$, the coefficients are
\[
 b_{k,m}=\sum_{j=0}^q a_j\frac{(k)_j}{(m)_j},
\]
where $(k)_j=k(k-1)\cdots(k-j+1)$ and the ratio is zero for $k<j$.
The ratio is the probability that $j$ draws without replacement from $m$
items all fall in a distinguished set of size $k$. The corresponding
with-replacement probability is $(k/m)^j$. Couple the two procedures until
the first repeated draw; its probability is at most $j(j-1)/(2m)$.
Hence uniformly in $k$,
\[
 |b_{k,m}-H(k/m)|\le
 \frac1{2m}\sum_{j=0}^q |a_j|j(j-1).
\]
Uniform coefficient bounds allow a common $m$ with this error less than
$h/2$. Thus $b_{k,m}\ge h/2$ and they are uniformly bounded above.

Accepted coefficients are multiplied by the positive scalar $g(x)$;
failure coefficients are the integrals of $(1-g)$ times the original
coefficients. Both operations preserve positive uniform bounds when the
reference measure is finite and the gate is bounded away from zero and one.
Multiplication of degree-$m$ Bernstein polynomials gives the same convolution
identity as \eqref{eq:bern-product}. Dividing by $D^n$ gives the exact
likelihood, and multiplication by a fixed prior followed by normalization
gives the posterior. There is no finite-rank or horizon-independent claim.
\end{proof}

In the cartridge model $D=a_0$ and the explicit degree $m=3$ works without
elevation. In the historically studied experiment conditional on successful
placement, $D=A(R)$ and the two remaining numerators are again cubic.
The same coefficient algorithm then retains the factor $A(R)^{-n}$ in the
posterior; $n$ counts prepared experiments rather than ambient placement
attempts. The two protocols are not equated by deleting this denominator.

\subsection{A matching local dimension lower bound}
The coefficient count is not merely an upper bound obtained by saving an
arbitrary collection of tests. The following worst-case lower bound applies
to continuous exact likelihood states for the entire command class.

\begin{theorem}[Sharp continuous state dimension at a fixed budget]\label{thm:dimension}
Assume the mode set contains a mode with $T>0$ and $\epsilon>0$. For each
$n\ge1$, normalized likelihoods attainable after $n$ observed censoring
failures contain a nonempty open subset of the degree-$3n$ coefficient
simplex. Thus they have local dimension $3n$. Any Euclidean state that
continuously encodes the known command history and reports, and exactly
determines the likelihood curve up to a scalar, needs at least $3n$ real
coordinates at that budget. The normalized $3n+1$ coefficient representation
has exactly $3n$ free coordinates and attains this bound. The same conclusion
holds for a posterior-determining state under a fixed full-support prior.
\end{theorem}
\begin{proof}
Fix $a$ with $b=2T>0$ and $\epsilon>0$. Consider the smooth gate family
\[
 g(\partial)=\tfrac12+\alpha_\partial,\quad
 g(0,y)=\tfrac12+\alpha_0,\quad
 g(1,y)=\tfrac12+\alpha_1+\alpha_2\cos(y-\theta).
\]
In a neighborhood of zero all gates satisfy the strict bounds. Their four
moments are $(1/2+\alpha_\partial,1/2+\alpha_0,1/2+\alpha_1,\alpha_2/2)$.
The map from these moments to the four power coefficients of
$H_\dagger$ is invertible: the constant coefficient determines $m_0$,
the linear coefficient then determines $m_1$, the quadratic coefficient
determines $m_\partial$, and the cubic coefficient determines $m_2$.
Consequently the attainable failure cubics contain an open neighborhood of
the constant polynomial $a_0/2$ in the four-dimensional polynomial space.

Choose $n$ polynomials $F_i(R)=a_0/2+\gamma_iR^3$ in this neighborhood,
with distinct small nonzero $\gamma_i$. They have degree three and are
pairwise coprime: a common root for two would force the two $\gamma_i$ to
coincide. The derivative of their product map is
\[
 (\delta F_1,\ldots,\delta F_n)\longmapsto
                 \sum_i\delta F_i\prod_{j\ne i}F_j.
\]
If this polynomial vanishes, reducing it modulo $F_i$ and using coprimality
shows that $F_i$ divides $\delta F_i$. Since both have degree at most three,
$\delta F_i=\lambda_iF_i$. The remaining condition is
$\sum_i\lambda_i=0$. The kernel has dimension $n-1$, so the product map
has rank $4n-(n-1)=3n+1$, the full dimension of polynomials of degree at most
$3n$. Normalization by the positive Bernstein coefficient sum removes one
scalar and has rank $3n$. The submersion theorem therefore gives an open
set of normalized positive coefficient vectors, with a continuous local
section back to their gate command histories. Each such all-failure history
has positive probability, at least $\eta^n$, at every radius.

An exact continuous state composed with that local section is a continuous
injection from an open subset of $\R^{3n}$ into its state space: equal states
would give equal normalized likelihoods. There is no continuous injection
of an open subset of $\R^{3n}$ into $\R^k$ for $k<3n$. Indeed, padding such
an injection by zeros would contradict invariance of domain in
$\R^{3n}$. A full-support prior also identifies continuous positive
polynomials up to a scalar from their normalized posterior densities:
equality almost everywhere extends to the interval by continuity and full
support. This proves the posterior version.
\end{proof}

The bound is across the command family, not a claim that the histories of
one fixed controller fill this open set. It excludes neither discontinuous
set-theoretic encodings into a single real number nor approximate state
reduction. It specifically establishes the sharp local continuous dimension
for the full experiment family and its known commands.
